\documentclass[11pt, t, handout]{beamer}
\usepackage{amsmath,amssymb,graphicx,tikz}
\usecolortheme{rose}
\useinnertheme[shadow]{rounded}
\usepackage[T1]{fontenc}
\usepackage[utf8]{inputenc}
\usepackage{booktabs}
\usepackage{xcolor}
\addtobeamertemplate{frametitle}{}{\vspace{0em}}
\setbeamertemplate{navigation symbols}{}
\makeatletter
\setbeamertemplate{footline}{%
  \ifnum\insertframenumber>1
    \begin{tikzpicture}[remember picture, overlay]
      \node[anchor=south, inner sep=4pt, yshift=3pt] at (current page.south) {%
        {\tiny\color{structure.fg}\insertframenumber\,/\,\inserttotalframenumber}};
    \end{tikzpicture}
  \fi
}
\makeatother
\newcommand{\ans}[1]{\textcolor{structure.fg}{#1}}
\newcommand{\todo}[1]{\textcolor{red!70!black}{\textbf{TODO:} #1}}

\title{A Risk Score for Depression Outcomes \\ \large Stage 1 of a PATH-guided two-stage IPD model}
\author{Alexander Richard}
\institute{DSSGx Munich, Collaborative Care IPD}
\date{\today}

\begin{document}

\begin{frame}\titlepage\end{frame}

% ─────────────────────────────────────────────────────────────────────────────
\begin{frame}{Where this fits}
  \pause
  \begin{center}\footnotesize
  \begin{tabular}{p{3.6cm}p{6.5cm}}
    \toprule
    \textbf{Stage 1 (this deck)} & Predict outcome risk from baseline
    characteristics only, no treatment information. \ans{The subject of
    everything that follows.} \\
    \midrule
    \textbf{Stage 2 (future work)} & Ask whether the risk score modifies
    the treatment effect; likely combined with a component network
    meta-regression across the collaborative-care building blocks. \\
    \bottomrule
  \end{tabular}
  \end{center}
  \vspace{10pt}
  \pause
  {\footnotesize Kent, Paulus, van Klaveren et al., \emph{The PATH
   Statement}, Ann Intern Med 2020: predict risk first, on data that
   cannot leak treatment assignment; only then examine effect
   modification.}
\end{frame}

% ─────────────────────────────────────────────────────────────────────────────
\begin{frame}{Why a two-stage IPD design, not one pooled regression}
  \pause
  \todo{State the between-study heterogeneity we actually observe (visit
  schedules, instrument families, predictor coverage) once the cohort for
  the chosen scenario is finalised.}
  \vspace{8pt}
  \pause
  {\footnotesize Debray, Moons, Ahmed, Koffijberg \& Riley, \emph{A
   framework for developing, implementing and evaluating clinical
   prediction models in an IPD meta-analysis}, Stat Med 2013; Riley,
   Tierney \& Stewart (eds), \emph{IPD Meta-Analysis: A Handbook for
   Healthcare Research}, Wiley 2021.}
\end{frame}

% ─────────────────────────────────────────────────────────────────────────────
\begin{frame}{Cohort definition: the scenario in force}
  \pause
  \todo{One table per finalised scenario: instrument, baseline/follow-up
  target times and tolerance, predictor list, missing-data policy,
  studies retained, N. Generated directly from cohort\$flow -- do not
  retype these numbers by hand.}
\end{frame}

% ─────────────────────────────────────────────────────────────────────────────
\begin{frame}{Assumptions this analysis makes}
  \pause
  \todo{Enumerate explicitly, e.g.: visits within the tolerance window are
  exchangeable with the nominal target time; a predictor missing for an
  entire study is missing at random across studies conditional on
  observed study-level characteristics (or: state why this is doubtful);
  complete-case analysis at Stage 1 does not introduce differential
  selection bias by outcome.}
\end{frame}

% ─────────────────────────────────────────────────────────────────────────────
\begin{frame}{Sample size}
  \pause
  \todo{Report the pmsampsize calculation once outcome type and target
  performance (shrinkage, expected $C$/$R^2$) are decided, against the
  realised N per scenario.}
  \vspace{8pt}
  \pause
  {\footnotesize Riley, Ensor, Snell et al., \emph{Calculating the sample
   size required for developing a clinical prediction model}, BMJ 2020;
   Riley, Snell, Ensor et al., \emph{Minimum sample size for developing a
   multivariable prediction model, Part II}, Stat Med 2019.}
\end{frame}

% ─────────────────────────────────────────────────────────────────────────────
\begin{frame}{Missing data}
  \pause
  \todo{State the Stage-1 policy (complete-case, per
  enforce\_non\_missingness) and its cost in dropped studies/patients;
  outline the Stage-2 plan if predictors turn out to be only
  sporadically missing rather than systematically absent.}
  \vspace{8pt}
  \pause
  {\footnotesize Jolani, Debray, Koffijberg \& Moons, \emph{Imputation of
   systematically missing predictors in an IPD meta-analysis}, Stat Med
   2015; Resche-Rigon \& White, \emph{MICE for systematically and
   sporadically missing multilevel data}, Stat Methods Med Res 2018.}
\end{frame}

% ─────────────────────────────────────────────────────────────────────────────
\begin{frame}{Model}
  \pause
  \todo{Specify the fitted model once the outcome definition is settled:
  multilevel structure (random intercept/slope by study), prior choices
  if Bayesian, shrinkage method if frequentist.}
  \vspace{8pt}
  \pause
  {\footnotesize Steyerberg, \emph{Clinical Prediction Models}, 2nd ed.,
   Springer 2019; B\"urkner, \emph{brms: An R Package for Bayesian
   Multilevel Models}, JSS 2017.}
\end{frame}

% ─────────────────────────────────────────────────────────────────────────────
\begin{frame}{Validation}
  \pause
  \todo{Report internal-external cross-validation (leave-one-study-out):
  per-study calibration and discrimination, then their meta-analytic
  summary and prediction interval.}
  \vspace{8pt}
  \pause
  {\footnotesize Steyerberg \& Harrell, \emph{Prediction models need
   appropriate internal, internal-external, and external validation},
   J Clin Epidemiol 2016; Riley et al., \emph{External validation of
   clinical prediction models using big datasets from e-health records or
   IPD meta-analysis}, BMJ 2016.}
\end{frame}

% ─────────────────────────────────────────────────────────────────────────────
\begin{frame}{Results}
  \pause
  \todo{Fill in once a scenario is fitted: coefficients/shrunk estimates,
  calibration plot, discrimination, IECV forest plot across studies.}
\end{frame}

% ─────────────────────────────────────────────────────────────────────────────
\begin{frame}{Limitations}
  \pause
  \todo{At minimum: unequal and partly continuous follow-up timing across
  trials; complete-case analysis under uneven predictor coverage;
  observational risk factors from RCT control-arm-heavy data; scale
  non-comparability across depression instruments left deliberately
  unbridged.}
\end{frame}

% ─────────────────────────────────────────────────────────────────────────────
\begin{frame}{Roadmap}
  \pause
  \begin{center}\footnotesize
  \begin{tabular}{ll}
    \toprule
    now & finalise outcome definition and candidate predictor set \\
    next & sample size target, then fit \\
    then & IECV, calibration, report \\
    stage 2 & risk-based effect modification + component NMA \\
    \bottomrule
  \end{tabular}
  \end{center}
\end{frame}

\end{document}
